\title{Controlling Text-to-Image Diffusion by Orthogonal Finetuning}
Recent years have witnessed significant advancements~\cite{saharia2022photorealistic,ramesh2022hierarchical,rombach2022high,gu2022vector,nichol2022glide} in generating images from text descriptions, largely propelled by developments in diffusion-based generative frameworks~\cite{ho2020denoising,song2021score,song2021denoising,dhariwal2021diffusion} and improved vision-language models~\cite{su2020vlbert,lu2019vilbert,radford2021learning,wang2021simvlm,li2022blip,li2023blip,alayrac2022flamingo,schuhmann2022laion}. For instance, LoRA requires 20 epochs to reach the level of accuracy attained by our OFT framework at only the 8th epoch. For each approach, we utilize the variant that obtains the highest validation CLIP score. Present approaches tend to operate under the unexamined premise that a finetuned model that remains close in Euclidean space to its pretrained counterpart is more likely to have preserved its original abilities, often leading practitioners to select a very small learning rate during finetuning. It is noteworthy that the number of trainable parameters in OFT and COFT is on par with LoRA—and much lower than that of ControlNet—while providing markedly better convergence efficiency compared to existing baselines. Within Transformers, for instance, LoRA is usually implemented by modifying the attention weights~\cite{hulora2022}. Further evidence of OFT’s rapid convergence is shown in Figure~\ref{fig:teaser}: the rightmost example demonstrates that OFT achieves strong convergence using just 5\% of the ADE20K dataset, highlighting its superior data efficiency over both ControlNet and LoRA. SphereNet~\cite{liu2017deep} demonstrates that when all neurons in a neural network are normalized to the unit hypersphere during training, the resulting network not only retains comparable representational capacity but also tends to generalize better. It is evident from the results that both COFT and OFT can stably adapt the diffusion model. In comparison, OFT maintains the integrity of the pretrained weight matrix by enforcing the transforming matrix to be orthogonal (and thus full-rank), ensuring that pre-existing information is preserved. To further enhance finetuning robustness, we introduce Constrained Orthogonal Finetuning (COFT), which incorporates an added radius restriction on the hypersphere. While this built-in regularization is intuitively sensible, it may remain overly permissive when adapting large models. Turning to the second question, a straightforward strategy for adjusting neuron directions is to collectively rotate or reflect all neurons in a particular layer, which amounts to an orthogonal transformation. First, whereas \cite{liu2021orthogonal} prioritizes reducing hyperspherical energy, OFT is designed to retain the pretrained model’s original hyperspherical energy, thereby protecting the pre-established semantic organization during finetuning; indiscriminately minimizing hyperspherical energy during diffusion model finetuning could potentially undermine the model’s learned semantics. For an orthogonal matrix with $r$ diagonal blocks, there are $\thickmuskip=2mu \medmuskip=2mu d(d/r-1)/2$ parameters, hence the parameter count scales as $\thickmuskip=2mu \medmuskip=2mu \mathcal{O}(d^2/r)$. Our results exhibit substantial gains over previous approaches in aspects of image quality, training stability, and speed of convergence. Overall, our observations indicate that LoRA exhibits a lower performance limit, even with extensive hyperparameter tuning. We adopt ControlNet~\cite{zhang2023adding}, T2I-Adapter~\cite{mou2023t2i}, and LoRA~\cite{hulora2022} as baseline approaches for comparison. Alternatively, OFT can be interpreted as modifying the canonical coordinate basis of layer neurons. \item For the controllable generation task, we propose a novel control consistency metric to quantify controllability. Drawing on prior empirical findings that hyperspherical similarity effectively represents semantic relationships~\cite{liu2017deep,liu2018decoupled,chen2020angular}, we leverage hyperspherical energy~\cite{liu2018learning} to quantify the relational structure between neurons within a layer. Our central conjecture is that an optimal finetuned model should display minimal change in hyperspherical energy when compared to its pretrained counterpart. By contrast, our focus is on adapting pretrained text-to-image diffusion models for improved controllability and enhanced generative outcomes in downstream tasks. In the following, we detail strategies to uphold orthogonality efficiently. \item We deploy OFT in the contexts of subject-driven generation and controllable generation, conducting an extensive empirical evaluation. Unlike frameworks such as ControlNet, the proposed OFT method does not introduce any inference-time overhead to the finetuned network. Conceptually, COFT incorporates two direct constraints: one limiting hyperspherical energy variation, and another explicitly restricting the deviation from the pretrained model’s parameters. Furthermore, OFT demonstrates improved sample efficiency, achieving strong performance with considerably fewer training images and iterations. In contrast, OFT places a different type of constraint focused on the similarity between pairs of neurons: $\thickmuskip=2mu \medmuskip=2mu \| \text{HE}(\bm{M})-\text{HE}(\bm{M}^0) \|=0$, with $\text{HE}(\cdot)$ denoting the hyperspherical energy of a weight matrix. Conversely, the absence of an orthogonality constraint in OFT results in an inability to synthesize plausible images or deliver controllable outcomes. While a straightforward strategy might involve regularizing hyperspherical energy to remain constant during finetuning, this does not guarantee that the actual energy difference is minimized. The main concept is to infer the control input from the generated image and directly compare it against the original control signal. Our approach is further influenced by orthogonal over-parameterized training~\cite{liu2021orthogonal}, which initiates classification neural networks from random orthogonal transformations and preserves the orthogonality throughout training. \textbf{Inference efficiency remains unaffected}. OFT builds upon the DreamBooth protocol but distinguishes itself by updating model parameters via orthogonal transformations instead of applying standard fine-tuning with a reduced learning rate. The baselines for comparison include DreamBooth~\cite{ruiz2023dreambooth} and LoRA~\cite{hulora2022}. \item In contrast to prevailing fine-tuning techniques, OFT modifies model parameters while theoretically guaranteeing the preservation of hyperspherical energy—a quantity we empirically observe to be indicative of maintaining the semantic generative capabilities inherited from the pretrained network. CLIP-T evaluates the average cosine similarity between the textual prompt and corresponding generated outputs via CLIP embeddings. This ensures the inference latency is equivalent to that of the original pretrained architecture. Notably, with an appropriate choice of $\epsilon$, COFT allows both the iteration count and learning rate to be selected with minimal effort, further enhancing usability. To provide a clearer illustration of the advantages offered by OFT, Figure~\ref{fig:dreambooth_final} presents a selection of randomly chosen subject-driven generation outputs. To assess convergence rates, we compare ControlNet, T2I-Adapter, LoRA, and COFT on the L2F task using both quantitative and qualitative analyses. To address this, we utilize an invariance property of hyperspherical energy: pairwise hyperspherical similarity remains unchanged when a shared orthogonal transformation is applied to all neurons. Figure~\ref{fig:eps_coft} illustrates how varying $\epsilon$ modulates the behavior of COFT: increasing $\epsilon$ enables the finetuned model to behave more like one adapted with OFT, while decreasing $\epsilon$ forces the finetuned model to remain closer to the original pretrained text-to-image diffusion model. The use of diffusion models has also been extended to this task~\cite{saharia2022palette,voynov2022sketch,wang2022pretraining}. \end{abstract}

\section{Introduction}

State-of-the-art text-to-image diffusion models~\cite{saharia2022photorealistic,ramesh2022hierarchical,rombach2022high} have established a new standard for producing high-quality images conditioned on textual inputs. \textbf{Controllable generation}. Our proposed forward computation is: $\thickmuskip=2mu \medmuskip=2mu\bm{z}=(\bm{R}\bm{W}^0\bm{D})^\top\bm{x}$\footnote{\textbf{Errata}: The NeurIPS camera-ready submission mistakenly presented the forward pass for re-scaled OFT as $\thickmuskip=2mu \medmuskip=2mu\bm{z}=(\bm{D}\bm{R}\bm{W}^0)^\top\bm{x}$. Firstly, orthogonality in OFT is ensured via Cayley parametrization, which entails computing a matrix inverse—a factor that modestly constrains its scalability. Taken together, these findings indicate that the OFT framework not only exhibits enhanced convergence behavior and stability, but also consistently delivers superior performance across evaluation metrics. Distinct from previous techniques, OFT guarantees the preservation of hyperspherical energy—a metric reflecting pairwise neuron interactions on the unit hypersphere. To further substantiate these findings, we include additional qualitative results for all control tasks in Appendix~\ref{app:control_exp}, and we further demonstrate the versatility of OFT by showcasing its application to additional control tasks (such as dense pose to human body, sketch to image, and depth to image) in Appendix~\ref{app:more_control_task}. Developing approaches that further improve parameter efficiency with reduced bias and increased adaptability constitutes an important open research direction. Figure~\ref{fig:toy_ae} demonstrates that altering neuron directions most effectively updates model semantics, which aligns with the design goals of OFT. Importantly, OFT incurs no extra inference costs during test-time execution. If we treat both $r$ and $r'$ as functions of the layer width $d$ (for instance, setting $\thickmuskip=2mu \medmuskip=2mu r = r' = \alpha d$ with some constant $\thickmuskip=2mu \medmuskip=2mu 0 < \alpha \leq 1$), the parameter complexity for LoRA is then $\mathcal{O}(d^2 + dn)$, whereas for OFT it scales as $\mathcal{O}(d)$. Although we alleviate this limitation through block diagonal parametrization, devising faster, differentiable methods for the matrix inverse remains unresolved. To mitigate this issue, we parameterize $\bm{R}$ as a block-diagonal matrix consisting of $r$ blocks, resulting in the following structure:

\begin{equation}
    \footnotesize
    \bm{R}=\text{diag}(\bm{R}_1,\bm{R}_2,\cdots,\bm{R}_r)=\begin{bmatrix}
    \bm{R}_1\in \text{O}(\frac{d}{r}) &  &\\
    &\ddots & \\
    & & \bm{R}_r\in \text{O}(\frac{d}{r})
    \end{bmatrix}\in \text{O}(d)
\end{equation}

Here, $\text{O}(d)$ refers to the set of $d$-dimensional orthogonal matrices. Models such as Stable Diffusion~\cite{rombach2022high} and DALL-E2~\cite{ramesh2022hierarchical} implement diffusion processes within a latent space. Moreover, the performance of OFT and COFT remains robust across a wide range of iteration counts, reducing the necessity for laborious iteration hyperparameter searches for different subjects. \textbf{Subject-driven generation}. For experimental clarity, we primarily employ the basic OFT setup to isolate the influence of orthogonal transformations alone, although our results indicate that the re-scaled variant generally outperforms the original (refer to Appendix~\ref{app:rescaled}). Our experimental results demonstrate that the OFT approach surpasses existing finetuning methods in terms of both image quality and convergence efficiency. Despite its strengths, OFT presents several compelling open challenges. Hyperparameters for DreamBooth and LoRA are set in accordance with recommendations from their original works, selecting the configurations with the best reported performance. LoRA outperforms ControlNet in the S2I task but is surpassed by it in both the C2I and L2F tasks. To promote fair comparisons, samples for each method are randomly drawn. The data show that our vanilla OFT maintains robust performance and precise control throughout training (notably at epoch 20). \textbf{Quantitative comparison}. Provided that $\bm{W}^0$ possesses full rank, OFT may also produce a low-rank adjustment if $\thickmuskip=2mu \medmuskip=2mu \bm{R}-\bm{I}$ is constrained to be low-rank. Our method specifically addresses two scenarios: subject-driven generation and controllable generation. Pivotal Tuning~\cite{roich2022pivotal} circumvents this by refining a generator near an inverted latent representation, with explicit regularization to maintain identity attributes. In particular, the orthogonal matrix is expressed as $\thickmuskip=2mu \medmuskip=2mu \bm{R}=(\bm{I}+\bm{Q})(I-\bm{Q})^{-1}$, where $\bm{Q}$ is a skew-symmetric matrix satisfying $\thickmuskip=2mu \medmuskip=2mu \bm{Q}=-\bm{Q}^\top$. Standard finetuning approaches offer maximal flexibility but at the cost of diminished stability and a greater risk of model collapse. All methods share the loss function specified by DreamBooth. It is important to note that cosine-based activations (option c) are not used during training, which relies solely on the inner product. Unlike the standard OFT in Eq.~\eqref{eq:oft_general}, where only $\bm{R}$ undergoes optimization, both $\bm{D}$ and $\bm{R}$ are now learnable parameters. To begin, we assess the stability and rate of convergence during finetuning for DreamBooth, LoRA, OFT, and COFT. Analogous to LoRA's practice of initializing with zeros, OFT must guarantee that the finetuning process originates precisely from the pretrained $\bm{W}^0$. In response, we present Orthogonal Finetuning~(OFT), a theoretically grounded finetuning strategy designed for customizing text-to-image diffusion models for specific tasks. Conversely, the OFT framework’s control effectiveness appears not to be capped, as empirical evidence suggests continued improvements with increased model capacity. As demonstrated in Table~\ref{table:control_gen}, OFT and COFT consistently deliver significantly better and more reliable control compared to other baseline methods. \subsection{Re-scaled Orthogonal Finetuning}

We introduce a straightforward modification to OFT by additionally optimizing a scaling factor per neuron. To ensure initialization of the orthogonal matrix $\bm{R}$ at the identity, $\bm{Q}$ is simply set to be the zero matrix initially. In line with \cite{ruiz2023dreambooth}, we further provide a quantitative analysis focusing on subject fidelity (DINO~\cite{caron2021emerging}, CLIP-I~\cite{radford2021learning}), prompt fidelity (CLIP-T~\cite{radford2021learning}), and output diversity (LPIPS~\cite{zhang2018unreasonable}). OFT diverges from the method in \cite{liu2021orthogonal} in two significant respects. To ensure robustness to stochasticity, each method undergoes 30 independent finetuning trials using varying random seeds, starting from the same pretrained checkpoint. For empirical evaluation, we select Stable Diffusion v1.5~\cite{rombach2022high} as the pretrained text-to-image backbone. Examining Figure~\ref{fig:dreambooth_final}, it becomes apparent that both OFT and COFT excel at retaining the subject's semantic features, whereas LoRA frequently struggles with maintaining subject identity (\emph{e.g.}, LoRA fails to preserve the subject in the bowl example). The phase spectrum—encoding the angular relationship between input signals and basis functions—primarily maintains the semantic information of the image. Traditional solutions largely rely on conditional generative adversarial networks~\cite{isola2017image,zhu2017toward,wang2018high,park2019semantic,choi2018stargan}. Within 400 iterations, DreamBooth and both OFT-based methods successfully exert meaningful control, whereas LoRA fails to maintain subject identity. Notably, the quantitative evaluation methodology for controllable generation presented here constitutes a significant novel contribution, providing an accurate means for measuring control fidelity in finetuned models—subject, of course, to the precision of the segmentation or detection models used for evaluation. We identify two main requirements for an effective finetuning strategy: (1) \emph{training efficiency}, characterized by minimizing both the number of trainable parameters and the duration of training in epochs; and (2) \emph{generalizability preservation}, ensuring that the model retains its capacity for producing high-quality, diverse outputs. In evaluation, we test three methods for constructing the feature map: (a) the standard inner product, as during training, (b) using only magnitude information, and (c) using only the angle information. Secondly, OFT exhibits a distinct capability for compositionality: the orthogonal matrices generated from different OFT finetuning tasks can be multiplied together, yielding another orthogonal matrix. This reliable and gradual convergence behavior of OFT significantly enhances its usability in practice, given that its training process is smoother and more predictable, making hyperparameter optimization considerably more straightforward. Preliminary results on using OFT for convolutional layer finetuning can be found in Appendix~\ref{app:convolution}. According to Table~\ref{table:dreambooth_final}, both COFT and OFT considerably surpass DreamBooth and LoRA with respect to DINO and CLIP-I metrics, and demonstrate either marginally superior or on-par results for prompt adherence and diversity. For the C2I task, evaluation is carried out using IoU and F1 scores, while the S2I task is measured by mean IoU as well as mean and overall accuracy. Figure~\ref{fig:face_convergence} presents the face landmark error as a function of the number of training epochs for each model. Likewise,~\cite{nitzan2022mystyle} develops a bespoke facial generator using a set of individual face images. Finetuning necessitates balancing adaptability with the stability of learned knowledge, and we contend that OFT offers a compelling solution to this challenge. Furthermore, \cite{liu2021orthogonal} establishes that orthogonal transformations possess sufficient expressive capacity for neural network training. Concretely, COFT uses the following computation during the forward pass:

\begin{equation}\label{eq:coft_general}
    \footnotesize
    \bm{z} = \bm{W}^\top \bm{x} = (\bm{R} \cdot \bm{W}^0)^\top \bm{x}, ~~~ \text{s.t.}~ \bm{R}^\top \bm{R} = \bm{R} \bm{R}^\top = \bm{I}, ~ \norm{\bm{R} - \bm{I}} \leq \epsilon
\end{equation}

Here, the transformation is subject to both the orthogonality constraint and an additional $\epsilon$-bound on deviation from the identity matrix. \subsection{Efficient Orthogonal Parameterization}\label{sect:eff_ortho}

Conventional orthogonalization procedures such as the differentiable Gram-Schmidt algorithm are frequently too computationally intensive for large-scale applications~\cite{liu2021orthogonal}. The underlying rationale for employing orthogonal transformation in neuron finetuning is partially informed by insights from the 2D Fourier transform, where an image is represented through its magnitude and phase spectra. Similarly, the T2I-Adapter~\cite{mou2023t2i} approach also finetunes diffusion models to allow for more precise control over image synthesis. To restrict this freedom, we advocate for preserving the angles between neuron pairs, drawing on the premise that these inter-neuronal angles are fundamental to a neural network's knowledge representation. Additional experimental results and information are available in Appendix~\ref{app:settings},\ref{app:coft_oft},\ref{app:more_qualitative},\ref{app:failure}. Additional information about these consistency metrics can be found in Appendix~\ref{app:settings}. Let $\thickmuskip=2mu \medmuskip=2mu \bm{R} \in \mathbb{R}^{d \times d}$ and, for each $i$, $\thickmuskip=2mu \medmuskip=2mu \bm{R}_i \in \mathbb{R}^{d/r \times d/r}$ denote orthogonal matrices. \textbf{Balancing flexibility and regularity in finetuning}. This observation underscores the importance of adjusting neuron orientations for semantic modifications in generated images, aligning with objectives in both subject-driven and controllable generation. To rigorously compare against LoRA, both OFT and COFT are only applied to the same network layers as LoRA. These qualitative findings highlight the superior balance achieved by OFT in terms of both controllability and subject fidelity. OFT, on the other hand, leverages a layer-shared orthogonal transformation to adjust neuron weights in a multiplicative way. In light of this, we propose learning a layer-wise, shared orthogonal transformation for neurons, ensuring that hyperspherical energy is maintained. \subsection{Controllable Generation}

\textbf{Settings}. \subsection{Subject-driven Generation}

\textbf{Experimental details}. While alternative angular modifications—such as rotating different neurons by distinct angles—are conceivable, we observe that orthogonal transformations effectively balance adaptability and structural discipline. However, these straightforward methods fall short of reliably maintaining the generative performance achieved by pretraining. Our principal contributions are outlined as follows:

\begin{itemize}[leftmargin=*,nosep]

    \item We introduce Orthogonal Finetuning, a new approach to fine-tune text-to-image diffusion models, enhancing their controllability. This insight is central to the OFT method: only an orthogonal transformation shared across

\begin{equation}\label{eq:oft_general}
    \footnotesize
    \bm{z}=\bm{W}^\top\bm{x}=(\bm{R}\cdot\bm{W}^0)^\top\bm{x},~~~\text{s.t.}~\bm{R}^\top\bm{R}=\bm{R}\bm{R}^\top=\bm{I}
\end{equation}

In this formulation, $\bm{W}$ represents the weight matrix updated through OFT, while $\bm{I}$ signifies the identity matrix. Each technique is employed to finetune Stable Diffusion~(SD) v1.5 on the respective datasets over 20 training epochs. This property has been supported by empirical findings in \cite{liu2021orthogonal}. Beyond diffusion-based systems, text-to-image synthesis has also involved generative adversarial networks~\cite{reed2016generative,zhang2017stackgan,xu2018attngan,li2019controllable} and autoregressive models~\cite{ramesh2021zero,ding2021cogview,wu2022nuwa,yuscaling}. Among finetuning techniques, adaptation methods (\emph{e.g.}, \cite{hulora2022,houlsby2019parameter,pfeiffer2020adapterfusion}) are well explored, particularly within natural language processing. In the LoRA approach with a low-rank dimension $r'$, the number of trainable parameters equals $\thickmuskip=2mu \medmuskip=2mu r'(d+n)$. \textbf{Model finetuning}. Despite these approaches aimed at improving efficiency, the matrix $\bm{R}$ retains its orthogonality, guaranteeing preservation of hyperspherical energy. For this layer, the output vector $\thickmuskip=2mu \medmuskip=2mu \bm{z}\in\mathbb{R}^n$ is computed by $\thickmuskip=2mu \medmuskip=2mu \bm{z}=\bm{W}^\top\bm{x}$ for input $\thickmuskip=2mu \medmuskip=2mu \bm{x}\in\mathbb{R}^d$. To assess controllable generation, we propose a control consistency metric, which quantifies how well the generated image reflects the original input control signal by extracting and comparing the control signal from the output image. An optional reduction in parameters is possible by tying the block matrices together, namely by enforcing $\thickmuskip=2mu \medmuskip=2mu\bm{R}_i = \bm{R}_j$ for all $i \neq j$, which leads to parameter complexity $\mathcal{O}(d^2/r^2)$. Our analysis reveals a fundamental tension between flexibility and regularity during finetuning. \item \textbf{Controllable generation}~\cite{zhang2023adding,mou2023t2i}: In this setting, an auxiliary control input—such as canny edge maps or segmentation layouts—guides the image generation process alongside the textual description. Additionally, OFT and COFT demonstrate substantially improved alignment with text prompt instructions, while DreamBooth, although effective at keeping subject identity, often cannot adhere to the content described by the text prompt (\emph{e.g.}, as shown in the first row of the bowl example). The strategy of adapting large pretrained models for specific downstream tasks has gained significant momentum~\cite{devlin2018bert,brown2020language,he2022masked}. Investigating whether this set of orthogonal matrices retains knowledge from each finetuned task presents a promising avenue for further study. While the orthogonality requirement can be satisfied using the Cayley parameterization (see Section~\ref{sect:eff_ortho}), imposing the $\epsilon$-deviation bound directly in the Cayley-parameterized formulation is more challenging. Conversely, sharing all blocks across $\bm{R}$ means every neuron undergoes a unified orthogonal transformation spanning channels. In contrast, ControlNet displays a delayed jump in controllability, only aligning with the landmarks after epoch 8—a phenomenon consistent with the sudden convergence described in \cite{zhang2023adding}. To gauge diversity, we compute the mean LPIPS-based cosine similarity among images of the same subject synthesized from identical prompts. For the quantitative evaluation, we calculate the mean $\ell_2$ distance between the predicted face landmarks and the reference control face landmarks. Of particular note, LoRA and OFT achieve parameter efficiency by fundamentally different means: LoRA leverages low-rank decomposition, whereas OFT employs a block-diagonal orthogonal structure, making use of inherent sparsity. In comparison to prior techniques, OFT achieves enhanced controllability and greater stability during finetuning, all while requiring fewer parameters. The corresponding outcomes are illustrated in Figure~\ref{fig:teaser} and Figure~\ref{fig:db_convergence}. Building on this property, we introduce Orthogonal Finetuning~(OFT), a method that enables the adaptation of large text-to-image diffusion models for specific tasks while maintaining invariant hyperspherical energy. This procedure rigorously maintains the pairwise angles between neurons in any given layer, contributing to enhanced stability. Next, we compare the parameter efficiency of OFT against LoRA. During inference, the learned orthogonal matrix $\bm{R}$ can be fused with the pretrained weights $\bm{W}^0$ via simple matrix multiplication, yielding $\thickmuskip=2mu \medmuskip=2mu \bm{W}=\bm{R}\bm{W}^0$. For the L2F task, our models achieve the highest accuracy in pose-controlled face generation, even under challenging orientations. Across these three control scenarios, both OFT and COFT outperform state-of-the-art baselines in qualitative image quality, attesting to the robustness of our OFT framework for controllable image synthesis. Specifically, hyperspherical energy is computed as the aggregate hyperspherical similarity (\emph{e.g.}, cosine similarity) across all neuron pairs in a layer, thereby reflecting how uniformly neurons are distributed on the unit hypersphere~\cite{liu2021learning}. Notably, OFT and COFT surpass ControlNet by producing images that exhibit richer detail and more plausible geometric coherence, all while using significantly fewer parameters. LoRA uses a rank parameter $r'$ to control parameter efficiency, while OFT instead utilizes a diagonal block parameter $r$ to restrict the number of parameters requiring training. We conduct a qualitative assessment of OFT and COFT against several leading approaches, such as ControlNet, T2I-Adapter, and LoRA. Within the OFT framework, the finetuning objective becomes minimizing the discrepancy in hyperspherical energy between the updated and original weights:

\begin{equation}\label{eq:oft_principle}
\footnotesize
\min_{\bm{W}} \left\| \text{HE}(\bm{W})-\text{HE}(\bm{W}^0)\right\|~~\Leftrightarrow~~\min_{\bm{W}} \bigg{\|} \sum_{i\neq j} \|\hat{\bm{w}}_i-\hat{\bm{w}}_j\|^{-1}-\sum_{i\neq j} \|\hat{\bm{w}}^0_i-\hat{\bm{w}}^0_j\|^{-1}\bigg{\|}
\end{equation}

where $\thickmuskip=2mu \medmuskip=2mu \hat{\bm{w}}_i:=\bm{w}_i/\|\bm{w}_i\|$ gives the normalized form of the $i$-th neuron, and the hyperspherical energy for $\bm{W}$ is given by $\thickmuskip=2mu \medmuskip=2mu \text{HE}(\bm{W}):= \sum_{i\neq j} \|\hat{\bm{w}}_i-\hat{\bm{w}}_j\|^{-1}$. In the special case when $\thickmuskip=2mu \medmuskip=2mu r = 1$, the resulting block-diagonal orthogonal matrix reduces to a full, unconstrained orthogonal matrix. To illustrate the centrality of neuron angles, we devised a simple experimental setup shown in Figure~\ref{fig:toy_ae}, employing a conventional convolutional autoencoder trained on a dataset of flower images. 
\begin{abstract}
Recent advances in large text-to-image diffusion models have led to remarkable success in synthesizing realistic images from textual descriptions. This extension is based on the observation that scaling neurons does not affect their hyperspherical energy, as any changes in magnitude are inherently normalized. However, allowing excessive freedom in altering neuron directions tends to compromise the generative capabilities established during pretraining. In addition, Appendix~\ref{app:human} describes a human evaluation study that provides additional empirical support for OFT's effectiveness. Examining the S2I task, it is apparent that LoRA fails to produce outputs in accordance with the provided segmentation maps. Furthermore, the field lacks a systematic framework for crafting robust finetuning techniques and for selecting crucial hyperparameters, including the appropriate number of training epochs. To better understand the structure of the Cayley transform, we can utilize the Neumann series to approximate $\thickmuskip=2mu \medmuskip=2mu \bm{R} =

As a consequence of the convergence of the series in operator norm, the constraint $\thickmuskip=2mu \medmuskip=2mu\|\bm{R}-\bm{I}\|\leq\epsilon$ can be reformulated within the Cayley transform, yielding the equivalent condition $\thickmuskip=2mu \medmuskip=2mu \|\bm{Q}-\bm{0}\|\leq\epsilon'$, where $\bm{0}$ indicates a matrix of all zeros and $\epsilon'$ serves as a distinct tolerance parameter from $\epsilon$. CLIP-I reflects the mean pairwise cosine similarity between CLIP embeddings for generated versus authentic images, while DINO operates similarly but substitutes ViT S/16 DINO representations. Nonetheless, these strategies frequently fail to retain the essential identity characteristics of the subject. Techniques based on inversion~\cite{choi2021ilvr,dhariwal2021diffusion,ramesh2022hierarchical,gal2022image} have been proposed to alter contextual details while keeping the main subject unaltered. Specifically, Stable Diffusion employs VQ-GAN~\cite{esser2021taming} to establish a latent codebook for visual representations, whereas DALL-E2 utilizes CLIP~\cite{radford2021learning} for constructing a shared embedding space linking images and text. For example, when each block matches an input channel, individual neuron channels are transformed separately, closely resembling the strategy of learning depth-wise convolutional filters~\cite{chollet2017xception}. Our method diverges from \cite{liu2021orthogonal} by not targeting minimization of hyperspherical energy. To ensure this, the orthogonal matrix $\bm{R}$ is initially set as an identity matrix, thereby maintaining that the finetuned parameters begin as the pretrained weights. \section{Concluding Remarks and Open Problems}

Inspired by the insight that the angular relationships among neurons are essential in shaping visual semantics, we introduce a straightforward yet powerful finetuning approach—orthogonal finetuning—for exerting control over text-to-image diffusion models. As a concrete illustration, consider a weight matrix of size $\thickmuskip=2mu \medmuskip=2mu 128 \times 128$; using LoRA with $\thickmuskip=2mu \medmuskip=2mu r' = 8$ results in $2,048$ tunable parameters, while OFT with $\thickmuskip=2mu \medmuskip=2mu r = 8$ (and no sharing among blocks) results in $960$ trainable parameters. Their findings demonstrate that orthogonal transformations provide ample capacity to train broadly generalizable classifiers. In contrast, the Orthogonal Finetuning (OFT) technique achieves an effective compromise: by maintaining the angles between neurons, OFT offers a simple yet robust balance between adaptability and regularization. This optimality is realized when $\bm{W}$ is a rotation or reflection of $\bm{W}^0$, i.e., $\thickmuskip=2mu \medmuskip=2mu \bm{W}=\bm{R}\bm{W}^0$, with $\thickmuskip=2mu \medmuskip=2mu \bm{R}\in\mathbb{R}^{d\times d}$ being an orthogonal matrix (where $\det(\bm{R})=1$ for a rotation and $-1$ for a reflection). \end{itemize}

At their core, both challenges revolve around the question of how to adapt text-to-image diffusion models through finetuning while maintaining the generative capabilities established during pretraining. Furthermore, OFT maintains efficiency by avoiding any increase in inference computational costs, making it highly suitable for practical deployment. To empirically substantiate this claim, we conduct a regularization experiment under the ControlNet~\cite{zhang2023adding} protocol, with results visualized in Figure~\ref{fig:ortho}. By the 2000-iteration mark, DreamBooth starts producing degraded or collapsed outputs, and LoRA misattributes yellow coloration to fur rather than clothing. To bridge this gap, this work addresses two key text-to-image generation challenges:

\begin{itemize}[leftmargin=*,nosep]

    \item \textbf{Subject-driven generation}~\cite{ruiz2023dreambooth}: Here, with access to only a small set of subject images, the model is tasked with creating novel renderings of the same subject in varying scenarios, utilizing textual prompts for context. Conversely, ControlNet, OFT, and COFT effectively align the generated images with the conditioning input. \subsection{Constrained Orthogonal Finetuning}

A stricter variant of OFT, termed COFT, further limits the flexibility of finetuning by ensuring that the adapted model remains within a fixed proximity to the original pretrained weights. Second, OFT emphasizes minimizing divergence from the pretrained weights, leading us to introduce a constrained formulation, while \cite{liu2021orthogonal} incorporates no analogous regularization. An overview of OFT is provided in Figure~\ref{fig:block}. For the qualitative comparison, our focus is on OFT, COFT, and ControlNet, as ControlNet records landmark errors closest to our results. Supplementary examples can be found in Appendix~\ref{app:more_qualitative}. \textbf{Why not minimize hyperspherical energy}. In sharp contrast, OFT and COFT remain capable of delivering reliable and steady control over image synthesis throughout. The most comparable to OFT is LoRA~\cite{hulora2022}, which implements additive low-rank updates to model weights during finetuning. We notice that adapter-based techniques (such as T2I-Adapter and ControlNet) tend to converge more slowly and ultimately achieve lower performance. To mitigate undesired modifications of specific subjects, methods such as~\cite{avrahami2022blended,nichol2022glide} integrate user-provided masks as auxiliary conditions. As illustrated in Figure~\ref{fig:control_final}, randomly sampled outputs indicate that OFT and COFT consistently deliver images with superior fidelity and realism while maintaining precise control over generation. This metric corroborates OFT's robust control capabilities. Ensuring the orthogonality of $\bm{R}$ can be accomplished using differentiable orthogonalization techniques as described by \cite{liu2021orthogonal,lezcano2019cheap}. This suggests that neuronal directions encode the most salient data characteristics. However, a key challenge remains: how can these models be effectively adapted or controlled to tackle diverse downstream applications? To improve efficiency, we leverage the Cayley parameterization for orthogonal matrix construction. \textbf{Qualitative comparison}. In practice, even though OFT achieves strong results, COFT demonstrates even greater stability during finetuning, which we attribute to this explicit deviation limitation. This method allows for efficient computation, but restricts the resulting orthogonal matrices to those with determinant $1$, i.e., matrices belonging to the special orthogonal group. Accordingly, adopting more structurally informed metrics for evaluating differences between pretrained and finetuned models holds promise for enhancing both the robustness and the quality of generative retention during finetuning. The approach in~\cite{hertz2022prompt} is capable of both global and localized edits without requiring masks as input. We focus our analysis on two representative finetuning scenarios: subject-driven generation—where the objective is to synthesize images depicting the same subject, given a handful of reference images and a text prompt; and controllable generation—which extends the model to process supplementary control signals in conjunction with textual guidance. While~\cite{casanova2021instance} enables instance variation synthesis, it may lose critical instance-level details. To address the first question, we draw upon empirical findings reported in \cite{liu2018decoupled,chen2020angular}, which indicate that the angular differences in features serve as a robust indicator of semantic disparity. In practice, this constraint does not compromise performance. \section{Orthogonal Finetuning}

\subsection{Why Does Orthogonal Transformation Make Sense?}

To motivate orthogonal transformation in the context of finetuning text-to-image diffusion models, we break the inquiry into two focused questions: (1) what is the rationale for finetuning the angular component of neurons (\emph{i.e.}, their directional properties), and (2) why is orthogonal transformation a suitable mechanism for altering these angles? These observations affirm that our OFT approach converges quickly while retaining stability in subject-conditioned generation tasks. This reformulation allows for straightforward enforcement of the $\bm{Q}$ constraint via projected gradient descent. To ensure a fair evaluation against LoRA, we restrict OFT to the finetuning of attention weights in our experiments. The results indicate that COFT and OFT both converge substantially more rapidly than the alternatives. As illustrated in Figure~\ref{fig:face_convergence_qual}, both OFT and COFT exhibit stable convergence, with the generated face pose incrementally matching the control landmarks over training. Essentially, the goal in conventional finetuning is to train the model while approximately minimizing $\thickmuskip=2mu \medmuskip=2mu \| \bm{M} - \bm{M}^0\|$, where $\bm{M}$ represents the finetuned weights and $\bm{M}^0$ the initial pretrained weights. Clearly, the objective in Eq.~\eqref{eq:oft_principle} reaches its global minimum of zero. DreamBooth~\cite{ruiz2023dreambooth}, on the other hand, uses a custom token and a limited set of subject images to adapt a text-to-image diffusion model, employing a combination of reconstruction and class-prior preservation losses. LoRA incorporates a low-rank matrix to ensure that the parameters of the pretrained weight matrix are not significantly altered during adaptation. The key approach involves learning a single orthogonal transformation per layer so as to keep neuron pairwise angles intact. \textbf{Why OFT converges faster}. These results confirm the critical role of the orthogonality constraint within the OFT framework. For L2F, we assess performance by computing the average $\ell_2$ distance between the generated landmarks and the input control landmarks. Despite the Cayley parameterization’s advantages, representing a full orthogonal matrix $\bm{R}$ remains parameter-inefficient when the dimension $d$ is large. Achieving minimal hyperspherical energy would distribute neurons uniformly on the hypersphere, which is not suitable for finetuning because it risks undermining the information obtained during pretraining. Lastly, the parameter efficiency achieved by OFT largely stems from the block diagonal parameterization, which inevitably introduces certain biases and restricts flexibility. We contend, however, that the Euclidean distance is an insufficient proxy for true semantic consistency. The primary paper evaluates three complex controllable generation scenarios: converting Canny edges to images~(C2I) using the COCO dataset~\cite{lin2014microsoft}, translating segmentation maps to images~(S2I) on the ADE20K dataset~\cite{zhou2017scene}, and generating faces from landmarks~(L2F) with the CelebA-HQ dataset~\cite{karras2018progressive,xia2021tedigan}. \textbf{Connection to LoRA}. GLIDE~\cite{nichol2022glide} and Imagen~\cite{saharia2022photorealistic} operate on the pixel domain, with GLIDE performing joint training of a text encoder and a diffusion process on paired text-image samples, while Imagen leverages a static, pretrained text encoder. Furthermore, OFT may be interpreted as conducting neuron optimization on a smooth hypersphere manifold, providing a more favorable optimization landscape. The total number of parameters for a $d \times d$ orthogonal matrix is $\thickmuskip=2mu \medmuskip=2mu d(d-1)/2$, leading to a complexity of $\mathcal{O}(d^2)$. Nevertheless, despite these advances, textual conditioning alone may not always yield precise or nuanced control over the generated content. These findings support the notion that the semantic representation of input images is predominantly embedded in the directional (angular) properties of neurons, suggesting that adjusting these directions is likely the most effective approach for altering semantic content. \textbf{Quantitative comparison}. Figure~\ref{fig:toy_ae} reveals that the angular information alone suffices to almost completely reconstruct the input images, whereas the magnitude information alone provides no meaningful content. Recently, ControlNet~\cite{zhang2023adding} has been proposed to enhance the controllability of pretrained diffusion models by finetuning them to accommodate additional control signals, leading to strong results. Beyond fully connected structures, OFT is also advantageous for adjusting convolutional layers, thanks to the meaningful role of the block-diagonal form of $\bm{R}$, which stands in contrast to the limitations of LoRA in this context. LoRA proposes a remedy by enforcing a low-rank structure on the parameter updates, specifically imposing $\thickmuskip=2mu \medmuskip=2mu \text{rank}(\bm{M} - \bm{M}^0)=r'$, where $r'$ is selected to be a small integer. \section{Experiments and Results}

\textbf{Experimental setup}. \end{itemize}

\section{Related Work}

\textbf{Text-to-image diffusion models}. Regarding the C2I task, both OFT and COFT generate semantically consistent images from coarse Canny edge maps, whereas T2I-Adapter and LoRA yield less satisfactory results. Finetuning is commonly approached by either applying a low learning rate when updating neural weights (\emph{e.g.}, \cite{ruiz2023dreambooth}) or by introducing compact modules that use re-parameterized weights (\emph{e.g.}, \cite{hulora2022,zhang2023adding}). Furthermore, OFT's performance remains robust to the number of iterations—large iteration counts do not negatively impact its results, unlike DreamBooth and LoRA, which both degrade under similar conditions. The original implementation is correct; thus, the findings in Appendix~\ref{app:rescaled} remain accurate.}, where $\thickmuskip=2mu \medmuskip=2mu\bm{D} = \text{diag}(s_1, \cdots, s_n) \in \mathbb{R}^{n \times n}$ is a trainable diagonal matrix with all entries $s_1,\cdots,s_n$ positive. This modification enhances OFT’s adaptability with a negligible increase in the parameter count. \section{Intriguing Insights and Discussions}

\textbf{OFT’s applicability across architectures}. In line with the objectives set by \cite{zhang2023adding,mou2023t2i}, our method applies OFT to finetune diffusion models, consistently delivering superior controllability, even with limited training data and fewer tunable parameters. Fundamentally, OFT does not rely on the specifics of any neural network architecture and can, in principle, be used with any model. \textbf{Convergence}. Although the latter constraint is commonly adopted implicitly by most finetuning strategies, COFT explicitly encodes it. Expanded results are reported in Appendix~\ref{app:more_qualitative}, \ref{app:more_control_task}, and \ref{app:failure}. The OFT framework is model-agnostic and readily applies to any text-to-image diffusion architecture. \end{document} For classification tasks, it is desirable for neurons to avoid redundancy. The experimental protocol for subject-driven generation follows that of DreamBooth~\cite{ruiz2023dreambooth}, whereas controllable generation settings align with those described in ControlNet~\cite{zhang2023adding} and T2I-Adapter~\cite{mou2023t2i}. Our investigation reveals that maintaining this property is essential to retain the original semantic image generation capabilities of diffusion models. \subsection{General Framework}

The standard finetuning approach generally employs gradient descent with a low learning rate to adjust the model parameters (or only select layers). We employ optimal hyperparameters for each method in comparison. Additionally, to reinforce training stability, we develop a constrained version that restricts the angular shift from the original pretrained weights. The OFT forward computation can be reframed as $\thickmuskip=2mu \medmuskip=2mu \bm{z}=(\bm{R}\bm{W}^0)^\top\bm{x}=(\bm{W}^0+(\bm{R}-\bm{I})\bm{W}^0)^\top\bm{x}$, where $\thickmuskip=2mu \medmuskip=2mu (\bm{R}-\bm{I})\bm{W}^0$ serves a similar purpose to the low-rank update used in LoRA. \textbf{Finetuning stability and convergence}. To illustrate this principle, consider a fully connected layer $\thickmuskip=2mu \medmuskip=2mu \bm{W}=\{\bm{w}_1,\cdots,\bm{w}_n\}\in\mathbb{R}^{d\times n}$, with each column vector $\bm{w}_i\in\mathbb{R}^{d}$ signifying the $i$-th neuron (and $\bm{W}^0$ indicating pretrained weights). \textbf{Qualitative comparison}. Random initialization is expected to yield low hyperspherical energy, and the central aim of \cite{liu2021orthogonal} is to maintain such low energy levels to foster better generalization performance in classification settings~\cite{liu2018learning,lin2020regularizing}. Notably, even if the magnitude spectrum is replaced by a random counterpart, retaining the original phase spectrum allows for the reconstruction of the image's semantic content. Notably, parameterizations based on block-diagonal structures further act as more restrictive forms of orthogonal matrix regularization. Employing a small learning rate serves as an implicit regularizer, restricting substantial departures from the pretrained parameters. Further implementation details and additional results can be found in Appendix~\ref{app:settings}. This difficulty is compounded by the absence of a clear metric to assess how well a model retains its pretrained generative strengths. During training, we compute the feature map using the standard inner product—where $z$ is the output from the convolution kernel $\bm{w}$ applied to input $\bm{x}$ within a sliding window. Furthermore, considering the insight that a reduced Euclidean distance between finetuned and pretrained models leads to greater retention of pretrained abilities, we propose a refined version—Constrained Orthogonal Finetuning~(COFT)—which restricts the finetuned model to a hypersphere of fixed radius centered at the pretrained neurons. In order to ensure fairness, all samples are derived from the same finetuned model across methods, and no individual text prompt optimization is performed for any specific output. Image-to-image translation serves as an example of controllable generation, where the output is steered based on input constraints.